\documentclass[journal]{IEEEtran}
\usepackage[left=0.3in,right=0.3in,top=0.3in,bottom=0.3in]{geometry}
\setlength{\columnsep}{10mm}
\usepackage{amsmath}
\usepackage{algorithmic}
\usepackage{array}
\usepackage{url}
\usepackage{hyperref}
\usepackage{float}
\usepackage{placeins}
\usepackage{graphicx}
\usepackage{listings}
\usepackage{color,xcolor}
\usepackage{tikz}


\hyphenation{op-tical net-works semi-conduc-tor}
\definecolor{dkgreen}{rgb}{0,0.6,0}
\definecolor{gray}{rgb}{0.5,0.5,0.5}
\definecolor{mauve}{rgb}{0.58,0,0.82}
\definecolor{bgcolor}{rgb}{0.97,0.97,0.9}


\lstset{frame=lrbt,
backgroundcolor=\color{bgcolor},
  aboveskip=3mm,
  belowskip=3mm,
  showstringspaces=false,
  columns=flexible,
  basicstyle={\small\ttfamily},
  numbers=none,
  numberstyle=\tiny\color{gray},
  keywordstyle=\color{blue},
  commentstyle=\color{dkgreen},
  stringstyle=\color{mauve},
  breaklines=true,
  breakatwhitespace=true
  tabsize=3,
  numbers=left,
  stepnumber=1,    
  firstnumber=1,
  numberfirstline=true
}

\newcommand{\ecefigure}[4]
{
	\begin{figure}[H]
	\centering 
	\includegraphics[width=#4\textwidth]{pic/#1}
	\caption{#2.\label{fig:#3}}
	\end{figure}
}

\begin{document}

\title{CEN415 Lab Report 2: Motors}
\author{Student 1 (ID), Student 2 (ID), Student 3 (ID), Teaching Assistant, Faculty}
\markboth{Abu Dhabi University}{Electrical and Computer Engineering Department}

\maketitle
    \begin{tikzpicture}[remember picture,overlay]
    \node[anchor=north west,yshift=-1pt,xshift=1pt]%
        at (current page.north west)
        {\includegraphics[height=20mm]{pic/logo.JPG}};
    \end{tikzpicture}
    
\begin{abstract}
In this abstract section, discuss the overall experience including a summary of the procedure you conducted in this lab, the results you obtained and analyzed, and the engineering conclusions you made. Also, summarize what have you learned and how does the lab related to the theory in the course.
\end{abstract}




\section{Summary of Conducted Experiment}
In this section, write about the experimental setup, the main equipment used, and the main connections made. Also, summarize the procedure followed in your own words indicating difficulties faced and how they were overcome. Include a hardware circuit diagram and a real-world image of the experimental setup. Refer to your hardware circuit using \autoref{fig:hw} and to your real/virtual circuit using \autoref{fig:cs}. Copy-paste and use the same code to include more images.


\ecefigure
{schematic.png}
{Hardware Circuit Diagram. Point out the main connections in the caption}
{hw}
{0.45}  


\ecefigure
{tinkercad.png}
{Connected Circuit on Tinkercad. Point out the main things to see in the caption}
{cs}
{0.4}  

For software-based labs, make sure you include your code in segments using Listing environments. You may use the code below to do so and copy-paste it to include more segments of your code:

\begin{lstlisting}[language=Octave]
t=0:0.01:10 ; 
f =1/2; 
x_t=4*cos(2*pi*f*t)+2*cos(6*pi*t);
subplot(211); plot(t,x_t);
title('Continuous Time Signal');
\end{lstlisting}

\section{Results and Data Analysis and Interpretation}
This section is devoted to your interpretation of the outcome of the experiment or project. The information from the data analysis is examined and explained. You should describe, analyze, and explain (not just restate) all your results. This section should answer the question 'what does the data tell me?' Describe any logical projections from the outcome, for instance, the need to repeat the experiments or to measure certain variables differently. Assess the quality and accuracy of your procedure. Compare your results with expected behavior and explain any unexpected behavior. 

Typical things done for analysis and interpretation of data are: 
\begin{itemize}
    \item Plotting/tabling expected/calculated outputs against measured/observed outputs.
    \item Confirming laws/formulas/facts by the comparison.
    \item Commenting on deviations and their reasons such as noise.
    \item Explaining/justifying the observations/phenomena/results based on applying engineering principles learned in the course.
    \item Responding to questions/prompts based on the observations from the experiments.
    \item If results are unexpected, explaining the debugging procedure followed to diagnose and fix the problem.
    \item Calculating statistics from the observed data to use it in driving conclusions such as the mean, the variance, ... etc.
    \item Fitting data to theoretical model or formula and explaining outliers.
    \item Commenting on the utility/limitations of a theoretical model/software/equipment in solving a problem.
    \item Discussing the sufficiency/insufficiency of testing to draw conclusions and identifying additional testing needs.
    \item Varying the values of the physical quantities in the system via the used software to investigate their impact on the system's output.
\end{itemize}


\subsection{First Exercise: (Replace with actual exercise title)}
In the first exercise, we observe... In \autoref{fig:Ohms} we can see...

We compare the experimental results with the expected theoretical results, and notice the following:
   \begin{itemize}
   \item The measured values are lower than the ideal values.
   \item The behaviour of the circuit is due to...
   \end{itemize}
   
Several errors could explain the differences: precision errors in measurement and bias errors in the calibration or large range of values considered. The data support/does not support the hypothesis because...

To add your table of results, you may use the \href{https://www.tablesgenerator.com/#}{latex tables generator website}. Once you create your table, you may set the options there then copy and paste your code here.

Observe an example of a table generated for voltage and current values for a fixed resistance in \autoref{tab:my-table}. The tabulated results have been plotted in Matlab as shown in \autoref{fig:Ohms}. The code for plotting your acquired results is given in \autoref{code}. You only need to fill your input and run. It will generate the plot and save it automatically to your directory. All styling options are set. \\
You may notice that the results support Ohm's Law given by \autoref{OhmsLaw} \cite{book1, book2}. We can observe that as the voltage increases, the current increases for the same resistance value \cite{paper1, paper2, paper3}. 

\begin{table}[!htbp]
\caption{Table caption goes here}
\label{tab:my-table}
\resizebox{\columnwidth}{!}{
\begin{tabular}{|c|c|c|c|c|c|}
\hline
\textbf{Voltage} & \textbf{Measured Voltage} & \textbf{Resistance} & \textbf{Measured Current} & \textbf{Calculated Current} & \textbf{Percentage Error} \\ \hline
3V & 3.01V & 1KOhm & 3.02mA & 3.01mA & 0.33\% \\ \hline
5V & 5.12V & 1KOhm & 5.2mA & 5.12mA & 1.5\% \\ \hline
7V & 7.05V & 1KOhm & 7.02mA & 7.05mA & 0.43\% \\ \hline
9V & 9.15V & 1KOhm & 9.3mA & 9.15mA & 1.6\% \\ \hline
\end{tabular}
}
\end{table}

\ecefigure
{ECE_figure.png}
{The current at different voltage levels. Give a sentence or two to help the reader focus their attention on what is the main interpretation from the plot (e.g., the relationship is observed as linear).}
{Ohms}
{0.45} 


\begin{equation}
\label{OhmsLaw}
    V = IR
\end{equation}




\section{Engineering Conclusions}
Base all conclusions on your actual results. Explain the meaning of the experiment and the implications of your results. Examine the outcome in light of the stated objectives. Seek to make conclusions in a broader context in the light of the results.

Throughout this lab report, three lab experiments were conducted in order to enhance our comprehension of this introductory course to electrical and computer engineering. The first lab along with the third lab experiments, acquaint us with... The second lab experiment emphasized on... In each lab, we set-up our connections to obtain... We achieve the following results, ..., which agree with the expected/theoretical values. We inspect several errors in measurements, which are due to... To correct problems that may have led to discrepancies or bad data points, we recommend the following... By conducting these lab experiments, we have learned..., which will help us with...

\appendices
\section{Source Code for Generating Matlab Figures Conforming to Required Style.}
\label{code}
\begin{lstlisting}[language=Octave]
clear; close all; clc
X = [3 5 7 9 12 15 18 21];
X_name  = 'X';
X_units = ' (X units)';
Y_name  = 'Y';
Y_units = ' (Y units)';
% Replace with your model
calculated = X.^2;
% Replace with your measurement 
measured = calculated+(20*rand(1,length(calculated))-10);
% mesured = [1 2 3 ...]

% Do not change here. This part is automatic. 
figure1 = figure;
% Create axes
axes1 = axes('Parent',figure1); hold(axes1,'on');
% Create multiple lines using matrix input to plot
plot1 = plot(X,[calculated; measured],'LineWidth',2,'Color',[0 0 0]);
set(plot1(1),'DisplayName','Calculated');
set(plot1(2),'DisplayName','Measured',
'Marker','o','LineStyle','--');
% Create labels
ylabel(strcat(Y_name,Y_units),'FontSize',20); xlabel(strcat(X_name,X_units),'FontSize',20);
% Create title
title(strcat(strcat(Y_name,' vs. '),X_name),'FontSize',20); box(axes1,'on');
% Set the remaining axes properties
set(axes1,'XGrid','on','YGrid','on');
% Create legend
legend1 = legend(axes1,'show');
set(legend1,...
    'Position',[0.290689149560117 0.795362903225807 0.172287390029325 0.0856854838709677],...
    'FontSize',16);

saveas(gcf,'ECE_figure.png')
\end{lstlisting}
\begin{thebibliography}{1}

\bibitem{IEEEhowto:kopka}
H.~Kopka and P.~W. Daly, \emph{A Guide to \LaTeX}, 3rd~ed.\hskip 1em plus
  0.5em minus 0.4em\relax Harlow, England: Addison-Wesley, 1999.
  
\bibitem{book1}
W. Hayt, J. Kemmerly and S. Durbin, Engineering circuit analysis, $11^{th}$ ed. Wiley, 2015.
  
\bibitem{book2}
R. Dorf and J. Svoboda, Introduction to electric circuits, $9^{th}$ ed. Hoboken, NJ: Wiley, 2014.

\bibitem{paper1}
M. Gupta, "Georg Simon Ohm and Ohm's Law", IEEE Transactions on Education, vol. 23, no. 3, pp. 156-162, 1980. 

\bibitem{paper2}
N. Kipnis, "A Law of Physics in the Classroom: The Case of Ohm’s Law", Science \& Education, vol. 18, no. 3-4, pp. 349-382, 2009. Available: 10.1007/s11191-008-9142-x 

\bibitem{paper3}
K. Griffiths, "Stray current control - an application of Ohm's law", IET Seminar on EMC in Railways, 2006. Available: 10.1049/ic:20060190

\end{thebibliography}

\end{document}


